\section{Математические модели механики жидкости и газа}
\label{sec:models-fluid-gas}

Механика жидкости и газа рассматривает движение и равновесие текучих
сплошных сред. Математическая модель строится из
\textbf{законов баланса массы, импульса и энергии}, определяющих
соотношений, уравнения состояния и начально-граничных условий.

Основные поля:
\[
\rho(\mathbf x,t),\qquad
\mathbf v(\mathbf x,t),\qquad
p(\mathbf x,t),\qquad
T(\mathbf x,t).
\]

\subsection{Гипотеза сплошной среды и кинематика}

В континуальном приближении молекулярная структура заменяется непрерывными
полями. Представительный физический объём должен быть мал по сравнению с
масштабом течения, но содержать большое число молекул. Критерием применимости
континуальной модели служит число Кнудсена
\[
Kn=\frac{\lambda}{L},
\]
где $\lambda$ --- длина свободного пробега, $L$ --- характерный размер. При
$Kn\ll1$ модель сплошной среды обычно применима.

Траектория жидкой частицы удовлетворяет
\[
\frac{d\mathbf x}{dt}=\mathbf v(\mathbf x,t).
\]
Для произвольного поля $\phi$ производная вдоль движения равна
\[
\boxed{
\frac{D\phi}{Dt}
=\frac{\partial\phi}{\partial t}
+(\mathbf v\cdot\nabla)\phi
}
\]
и называется \textbf{материальной производной}.

Градиент скорости раскладывается как
\[
\nabla\mathbf v=\mathbf D+\mathbf W,
\]
где
\[
\mathbf D=\frac12\left(\nabla\mathbf v+(\nabla\mathbf v)^T\right)
\]
--- тензор скоростей деформации, а антисимметричная часть $\mathbf W$
характеризует локальное вращение. Вихрь
\[
\boldsymbol\omega=\nabla\times\mathbf v.
\]

\subsection{Закон сохранения массы}

Для неподвижного контрольного объёма $V$
\[
\frac{d}{dt}\int_V\rho\,dV
+\int_{\partial V}\rho\mathbf v\cdot\mathbf n\,dS=0.
\]
С помощью формулы Гаусса получаем локальное уравнение неразрывности
\[
\boxed{
\frac{\partial\rho}{\partial t}
+\nabla\cdot(\rho\mathbf v)=0
}
\]
или
\[
\frac{D\rho}{Dt}+\rho\nabla\cdot\mathbf v=0.
\]
Для жидкости постоянной плотности
\[
\boxed{\nabla\cdot\mathbf v=0}.
\]

\subsection{Закон сохранения импульса}

Локальный баланс количества движения имеет вид
\[
\boxed{
\rho\frac{D\mathbf v}{Dt}
=\nabla\cdot\boldsymbol\sigma+\rho\mathbf f,
}
\]
где $\boldsymbol\sigma$ --- тензор напряжений, $\mathbf f$ --- массовая
сила на единицу массы. Для жидкости удобно выделять давление:
\[
\boldsymbol\sigma=-p\mathbf I+\boldsymbol\tau,
\]
поэтому
\[
\rho\frac{D\mathbf v}{Dt}
=-\nabla p+\nabla\cdot\boldsymbol\tau+\rho\mathbf f.
\]

\subsection{Закон сохранения энергии}

Пусть $e$ --- удельная внутренняя энергия, а
\[
E=e+\frac{|\mathbf v|^2}{2}
\]
--- полная удельная энергия. Уравнение внутренней энергии можно записать как
\[
\boxed{
\rho\frac{De}{Dt}
=-p\,\nabla\cdot\mathbf v
+\boldsymbol\tau:\mathbf D
-\nabla\cdot\mathbf q
+\rho r,
}
\]
где $\mathbf q$ --- тепловой поток, $r$ --- объёмный источник тепла на
единицу массы. Эквивалентная консервативная форма полной энергии:
\[
\frac{\partial(\rho E)}{\partial t}
+\nabla\cdot\bigl[(\rho E+p)\mathbf v\bigr]
=\nabla\cdot(\boldsymbol\tau\mathbf v)
-\nabla\cdot\mathbf q
+\rho\mathbf f\cdot\mathbf v+\rho r.
\]

\subsection{Определяющие соотношения и замыкание}

Законы сохранения универсальны, но для получения замкнутой системы нужны
свойства конкретной среды. Для ньютоновской жидкости
\[
\boxed{
\boldsymbol\tau
=2\mu\mathbf D
+\lambda_v(\nabla\cdot\mathbf v)\mathbf I,
}
\]
где $\mu$ и $\lambda_v$ --- коэффициенты вязкости. Для несжимаемой
ньютоновской жидкости
\[
\boldsymbol\tau=2\mu\mathbf D.
\]
Закон Фурье:
\[
\boxed{\mathbf q=-k\nabla T}.
\]
Для идеального газа
\[
\boxed{p=\rho RT},
\qquad
e=c_vT.
\]
Именно выбор определяющих соотношений отличает ньютоновские и
неньютоновские жидкости, идеальный газ и более сложные среды.

\subsection{Уравнения Эйлера и Навье--Стокса}

Если вязкостью и теплопроводностью пренебречь, получаются уравнения Эйлера.
Для сжимаемого потока:
\[
\frac{\partial\rho}{\partial t}+\nabla\cdot(\rho\mathbf v)=0,
\]
\[
\frac{\partial(\rho\mathbf v)}{\partial t}
+\nabla\cdot(\rho\mathbf v\otimes\mathbf v+p\mathbf I)
=\rho\mathbf f.
\]
Уравнение энергии дополняет систему.

Для несжимаемой ньютоновской жидкости постоянной плотности система
Навье--Стокса принимает вид
\[
\boxed{
\begin{aligned}
\nabla\cdot\mathbf v&=0,\\
\frac{\partial\mathbf v}{\partial t}
+(\mathbf v\cdot\nabla)\mathbf v
&=-\frac1\rho\nabla p+\nu\Delta\mathbf v+\mathbf f,
\end{aligned}
}
\]
где $\nu=\mu/\rho$ --- кинематическая вязкость. Нелинейность связана с
конвективным членом $(\mathbf v\cdot\nabla)\mathbf v$.

Для стационарного невязкого течения жидкости постоянной плотности под
действием консервативной массовой силы $\mathbf f=-\nabla\Pi$ справедлива
\textbf{теорема Бернулли}:
\[
\boxed{
\frac{p}{\rho}+\frac{|\mathbf v|^2}{2}+\Pi=\text{const}
}
\]
вдоль линии тока. В поле тяжести при вертикальной координате $z$, направленной
вверх, $\Pi=gz$, поэтому
\[
\frac{p}{\rho}+\frac{|\mathbf v|^2}{2}+gz=\text{const}.
\]
При безвихревом течении константа одинакова во всей связной области.

\subsection{Начальные и граничные условия}

Для нестационарной задачи задают начальные поля, например
\[
\mathbf v(\mathbf x,0)=\mathbf v_0(\mathbf x),
\qquad
\rho(\mathbf x,0)=\rho_0(\mathbf x).
\]
На неподвижной твёрдой стенке для вязкой жидкости используется условие
прилипания
\[
\mathbf v=0.
\]
Для идеальной жидкости обычно достаточно непротекания
\[
\mathbf v\cdot\mathbf n=0.
\]
На входе и выходе задают совместимый набор скоростей, расходов, давления,
температуры и других характеристик. Для сжимаемых течений число допустимых
граничных условий связано с направлениями характеристик.

\subsection{Безразмерные критерии}

После введения характерных масштабов $L$, $U$, $\rho_0$ у несжимаемых
уравнений Навье--Стокса появляется число Рейнольдса
\[
\boxed{Re=\frac{\rho_0UL}{\mu}=\frac{UL}{\nu}},
\]
характеризующее отношение инерционных и вязких эффектов. Для сжимаемых
потоков важно число Маха
\[
Ma=\frac{U}{a},
\]
для течений со свободной поверхностью --- число Фруда, а при совместном
переносе импульса и тепла --- числа Прандтля и Пекле. Безразмерные числа
определяют режимы течения и позволяют сравнивать физически разные системы.

\subsection{Турбулентность и численное моделирование}

При больших числах Рейнольдса движение может стать турбулентным. Основные
уровни моделирования:
\begin{itemize}
    \item DNS --- непосредственное разрешение всех существенных масштабов;
    \item LES --- разрешение крупных вихрей с моделью подсеточных масштабов;
    \item RANS --- осреднённые уравнения с моделью турбулентных напряжений.
\end{itemize}
Это пример \textbf{проблемы замыкания}: после осреднения возникают новые
неизвестные корреляции, для которых нужны дополнительные модели.

Численно применяются методы конечных разностей, конечных объёмов и конечных
элементов. Для законов сохранения особенно важна консервативность
дискретизации; для задач с сильной конвекцией --- устойчивость и корректное
описание разрывов.

\subsection{Что важно сказать на экзамене}

Ответ удобно строить как цепочку
\[
\boxed{
\text{сплошная среда}
\to
\text{масса, импульс, энергия}
\to
\text{определяющие соотношения}
\to
\text{граничные условия}
}.
\]
Обязательно записать уравнение неразрывности и систему Навье--Стокса для
несжимаемой жидкости, объяснить роль давления и вязкости. Полезно привести
теорему Бернулли с её условиями применимости и число Рейнольдса.

\newpage
